% !TITLE 相思
% !AUTHOR 王维
% !DYNASTY 唐
% !GENRE 五言绝句
% !ORDER 37

\begin{poem}
红豆\textsuperscript{1}生南国，春来发几枝。
愿君多采撷\textsuperscript{2}，此物最相思。
\end{poem}

\begin{annotations}
\notetext{1}{红豆：红豆树的种子，又名"相思子"，色红，古代用作爱情和友谊的信物。传说古时有人死于边塞，妻子在树下哭死化为红豆。}
\notetext{2}{采撷：采摘。撷（xié），摘下。希望你多采一些红豆，看到它就会想起我。}
\end{annotations}

\begin{appreciation}
这首诗据说是王维写给李龟年的。李龟年是唐代著名的音乐家，安史之乱后流落江南，王维以此诗寄赠。但诗一旦写成，便不再属于作者。千百年来，无数人用这首诗表达思念，它早已超越了具体的对象，成为相思本身的代名词。

红豆生南国，春来发几枝——红豆生长在南方，春天来了，又发出几枝新芽。红豆又名相思子，相传古时有人死于边塞，他的妻子在树下哭泣而死，化为红豆。所以红豆从诞生之日起，就带着悲伤的基因。发几枝，是一个问句，也是一个叹息：春天又来了，你还好吗？

愿君多采撷，此物最相思——希望你多采摘一些，这东西最能代表相思。采撷（cǎi xié）就是采摘。多采撷，不是要你吃，是要你看见。王维不说我想你，他说这豆子最相思。把情感投射到一个物件上，让物件替自己说话，这是中国诗歌最含蓄也最动人的方式。

这首诗只有二十个字，却写尽了人类最普遍的情感。它不依赖任何背景知识，不需要知道李龟年是谁，不需要知道红豆的传说。你只要知道：有一种红色的豆子，长在南方，春天发芽，看到它的人会想起远方的人。这就够了。

王维的诗常常如此：用最简单的语言，说最深沉的情感。他不渲染，不铺陈，只是轻轻一点，便让读者自己去填补那无限的空白。春来发几枝，问的是红豆，念的是人；愿君多采撷，说的是豆子，藏的是心事。这种含蓄，是中国诗歌独有的美。
\end{appreciation}
